\documentclass[11pt]{article}

\usepackage{fullpage}
\usepackage{amsmath}
\usepackage{epsfig}
\usepackage{graphics}
\usepackage{caption}
\usepackage{circuitikz}

\oddsidemargin -0.15in
\evensidemargin -0.15in
\headheight 0in
\headsep 0in
%\topmargin -0.25in
\textheight 9.3in
\textwidth 6.8in

\usepackage{enumitem,amssymb}
\newlist{todolist}{itemize}{2}
\setlist[todolist]{label=$\square$}

\begin{document}


\begin{center}
\section*{
Franklin \& Marshall College \hfill PHY 223\\
\bigskip
{\Large \it BB Scattering: Writing  (Part 1)}}
\end{center}

Today, you will start working on your lab report. Make sure to refer to the \textit{Guide to writing a lab report} as you do so, especially the section on ``Data and analysis'' and the style guide for figures.

\vspace{3mm}

\noindent {\Large {\bf Before starting to write}}

\vspace{3mm}

Your first task today is to take stock of where you are at. Go through the checklist of things needed before you start writing below:

\vspace{3mm}

\begin{todolist}

\item Finish taking your data.

\item Complete your analysis of the results, including full consideration of uncertainty in your data and derived parameters.

\item Create an initial set of figures for your report. These will include figures that show your results as well as schematic diagrams showing your experimental setup.

\item Based on your results, write down your claim. This should
  be the answer to a well-formed question. When possible, your claim
  should include numbers.

\item Revise your figures based on your claim. Do the figures present the
  right supporting information? Are all the figures relevant to the narrative?

  \item Think about your target audience. Consider the level of knowledge
 of the people you are writing for.

\item Review the literature in order to provide context and motivation for your work.

\end{todolist}

% \vspace{3mm}

{\it We will save reviewing the literature for the end of today, if time allows, but you should try to finish all the other bullet points.}

% \newpage
\vspace{6mm}

\noindent {\Large {\bf Making figures}}

\vspace{3mm}

Starting your report with the figures will help you think about what you need to tell your readers. It will center your thinking. Each of your reports should have a minimum of 3 figures:
\begin{itemize}
 \item At least one figure showing your data and results of your analysis.
 \item An experimental setup schematic.
 \item A figure that helps you explain the physics.
\end{itemize}

The number of figures in a report will vary greatly depending on the nature of the work. Even in this report, you could have more than three figures (though you want to avoid superfluous graphs). However, the three proposed here are a good starting point.

Create those three figures. The data and analysis figure(s) should be done using Python. The other two can often be done using Inkscape, PowerPoint, etc. For this report, make sure to \textbf{use Inkscape} to generate at least one of your figures in a vector-based format (e.g., SVG or PDF), although in some cases using a pixel-based software (e.g., Canva, Photoshop) can also be appropriate. \textit{You cannot hand-draw your figures}, but you can (and probably should) sketch your ideas first and discuss them with your team. Computer drawing can be time-consuming. A little bit of planning can save you a lot of time. If you have ideas for a fourth figure, discuss it with the instructor first.

\vspace{3mm}

As you make your figures, be sure to consider the practices we have discussed for making effective figures, most notably \textit{your audience}, which you should assume is a group of students with similar background knowledge to our class, but who have not completed the BB Scattering experiment themselves. Also consider the following questions:
\begin{itemize}
 \item Is the figure free of errors?
 \item Is the figure visually appealing?
 \item Are there elements of the figure that could be misleading or lead to misunderstandings?
 \item Are there colors or font sizes that are hard to read or interpret?
 \item Are the axes, ticks, grids, legends, etc. useful for communicating the results, or are they distracting?
 \item How will your figure help make a point in the report? \textit{Note: you may wish to revisit this point after you determine your Claim in the next section.}
\end{itemize}


Place your figures in a \textbf{single document} and write a \textbf{short caption} for each of them. The caption does not need to be the final version of the caption for your report, but it should be sufficient to explain what is going on in the figure. You will upload this document later today.

\vspace{5mm}

\noindent {\Large {\bf Claim}}

\vspace{3mm}

Write down your claim for this experiment.  \textit{Note: in this case, your claim should include
  a comparison to the measured size of the scattering object.} Talk to your instructor about ways of making this measurement, and make sure your claim makes sense given your results (including uncertainties).
  
  \textbf{Add your written claim to the document with your figures above.} Then submit your document as a PDF on Canvas in the appropriate assignment.

\vspace{5mm}

\noindent {\Large {\bf Background reading}}

\vspace{3mm}

Before you start writing, it is a good idea to have done some background reading to put your work into context. For this experiment, you will likely not find a whole lot about the Physics of BB scattering. It is a demonstration system meant to introduce more general concepts about scattering. But that doesn't mean we can't discuss the Physics of what is happening. There are many different types of scattering experiments. In all cases, we have a source of particles (or waves), a target we want to learn about, and a theory that allow us to interpret the results. You can look into this by searching online for:

\begin{itemize}
 \item Scattering
 \item Neutron scattering
 \item Rutherford experiment
 \item Etc.
\end{itemize}

None of these processes are completely applicable to what we are doing, and the theory behind each is very different. However, knowing a little about them will be helpful in situating your work and motivating why this kind of experiment is important. Take notes on these experiments and what you might want to include about them in your report to help contextualize and motivate your experiment.

\vspace{5mm}

% \noindent {\Large {\bf Editing figures}}

% Now that you have received comments on the figures and have written down your claim, make edits to the figures as necessary.

% {\bf Name your files, analysis$\_$final.png, exp$\_$setup$\_$final.png, and theory$\_$fig$\_$final.png. Then upload them to Canvas in the appropriate submission.}

\end{document}
